\documentclass[11pt]{article}
\usepackage[margin=1in]{geometry}
\usepackage{amsmath,amssymb,amsfonts}
\usepackage{bm}
\usepackage{amsthm}
\theoremstyle{remark}
\newtheorem{remark}{Remark}[section]
\usepackage{enumitem}
\usepackage{xurl}
\usepackage[colorlinks=true,linkcolor=blue,citecolor=blue,urlcolor=blue]{hyperref}

% file and column names: typewriter, breakable at any character (xurl)
\newcommand{\src}[1]{\path{#1}}
% a file of Section 1, by its label: F<number>
\newcommand{\fref}[1]{\ref{#1}}

\title{The age- and location-stratified SEIR model: the data\thanks{Parts of this
manuscript were edited with the assistance of an AI tool (Claude Opus 5.5, Anthropic).
The authors are solely responsible for all proofs and statements.}}
\author{Zhilan \and John \and Yuliya \and Ludmil}
\date{}

\begin{document}
\maketitle

\begin{abstract}
Description of the process of collecting public data and generating mock data.
\end{abstract}

\section{Data}\label{sec-data}
There are $m=51$ locations (the 50 states and D.C., in the order of their FIPS codes)
and $n=16$ age groups ($0$--$4$, \dots, $70$--$74$, $75+$).

\subsection{Files}\label{sec-files}

\subsubsection{Downloaded}\label{sec-downloaded}
Kept as downloaded, in
\src{MathBio/data/Downloaded_from_source/}; paths relative to this directory.
\begin{enumerate}[label=F1-\arabic*, ref=F1-\arabic*, series=files, leftmargin=*]
\item\label{f-jtb} \src{papers/2026FengDayGlasser.pdf}: the model paper~\cite{2026Feng_Day_Glasser}.
\item\label{f-jtbs} \src{papers/2026FengDayGlasserSupplement.pdf}: its supplement (Tables S1--S4).
Sections~\ref{sec-age}, \ref{sec-const}, \ref{sec-var}.
\item\label{f-g23} \src{papers/2023GlasserFengVoJonesClarke.pdf}:~\cite{2023Glasser_Feng}; its
Fig.~4 is digitized for $\beta^0$. Sections~\ref{sec-time}--\ref{sec-mock}.
\item\label{f-f17} \src{papers/2017FengHillCurnsGlasser.pdf}:~\cite{2017Feng_Hill_Curns_Glasser}.
Section~\ref{sec-loc}.
\item\label{f-prem} \src{papers/2017PremCookJit.pdf}:~\cite{2017Prem_Cook_Jit}, and its
supporting text S1 \src{papers/2017PremCookJitSupplementS1.pdf}. Section~\ref{sec-contacts}.
\item\label{f-baj}
\src{papers/2021BajemaWiegandCuffePatelIachanLimLeeMoyseHaversHardingFryHallMartinBielDengMeyerMathurKyleGundlapalliThornburgPetersenEdens.pdf}:~\cite{2021Bajema_Wiegand},
and its online supplement (eMethods; eTables 3--6, the seroprevalence by
jurisdiction and age group in the four periods)
\src{papers/2021BajemaWiegandCuffePatelIachanLimLeeMoyseHaversHardingFryHallMartinBielDengMeyerMathurKyleGundlapalliThornburgPetersenEdensSupplement.pdf}.
Section~\ref{sec-mock}.
\item\label{f-kingman} \src{papers/1961Kingman.pdf}:~\cite{1961Kingman-aa}, the
log-convexity of the spectral radius used in \src{stratified_seir_ngm.tex}
(Lemma~5.1); not a data source.
\item\label{f-sc} \src{census/sc-est2024-alldata6.csv.gz}: state population estimates by age,
Vintage 2024~\cite{2024Census_scest2024}. Section~\ref{sec-pop}.
\item\label{f-co} \src{census/co-est2024-alldata.csv.gz}: county population estimates,
Vintage 2024~\cite{2024Census_coest2024}. Section~\ref{sec-loc}.
\item\label{f-gaz} \src{census/2024_Gaz_counties_national.zip}: county land areas,
2024 Gazetteer~\cite{2024Census_gazetteer}. Section~\ref{sec-loc}.
\item\label{f-cenpop} \src{census/CenPop2010_Mean_ST.txt.gz}: 2010 centers of population of
the states~\cite{2010Census_cenpop}. Section~\ref{sec-loc}.
\item\label{f-premx} \src{prem2017/pcbi.1005697.s002.zip}: supporting file S2
of~\cite{2017Prem_Cook_Jit}, the contact matrices; the one used is
\src{MUestimates_all_locations_2.xlsx} in it. Section~\ref{sec-contacts}.
\item\label{f-ox} \src{oxcgrt/OxCGRT_compact_subnational_v1.csv.gz}: the final OxCGRT data
set~\cite{2021Hale_OxCGRT,2023OxCGRT_data}. Section~\ref{sec-loc}.
\item\label{f-sero} \src{cdc/seroprevalence.csv.gz}: the CDC Nationwide Commercial Laboratory
Seroprevalence Survey~\cite{2023CDC_sero}. Section~\ref{sec-mock}.
\item\label{f-var} \src{cdc/variants_usa.csv.gz}: CDC SARS-CoV-2 Variant
Proportions, USA, weighted estimates~\cite{2023CDC_variants}. Section~\ref{sec-var}.
\end{enumerate}
The data files are kept as downloaded, the text files gzipped (\src{.gz}) and the
zip archives as they are; the scripts read them without unpacking.

\subsubsection{Built}\label{sec-built}
The inputs of the model, in \src{MathBio/code/python/data/},
each written by the script named, from the files listed.
\begin{enumerate}[label=F1-\arabic*, ref=F1-\arabic*, resume=files, leftmargin=*]
\item\label{f-N} \src{census/N_51x16.csv}: $N_{ij}$; \src{census_N_51x16.py}
from~\fref{f-sc}. Section~\ref{sec-pop}.
\item\label{f-U} \src{prem2017/U_16x16.csv}: $U$; \src{prem_U_16x16.py}
from~\fref{f-premx}. Section~\ref{sec-contacts}.
\item\label{f-d} \src{census/d_51x51.csv}: $d_{ip}$; \src{geo_51.py}
from~\fref{f-cenpop}. Section~\ref{sec-loc}.
\item\label{f-L} \src{census/L_51.csv}: $\rho_i$ and $L_i$; \src{geo_51.py}
from~\fref{f-co} and~\fref{f-gaz}. Section~\ref{sec-loc}.
\item\label{f-C} \src{model/C_816x816.csv.gz}: the contact matrix $C$ (gzipped CSV; rows
$(i,j)$ and columns $(p,q)$ labelled by state and age group); \src{build_C_51x16.py}
from~\fref{f-U}, \fref{f-d}, \fref{f-L}. Section~\ref{sec-loc}.
\item\label{f-s} \src{oxcgrt/s_51_daily.csv}: $s_i(t)$, daily; \src{oxcgrt_s_51.py}
from~\fref{f-ox}. Section~\ref{sec-loc}.
\item\label{f-t2} \src{glasser2023/table2.csv}: $p_j$, $d_j$, $\mathrm{IFR}_j$, $c_j^0$, typed
from Table~2 of~\fref{f-g23}; \src{glasser2023_inputs.py}. Section~\ref{sec-age}.
\item\label{f-beta} \src{glasser2023/beta0.csv}: $\beta_j^0$, digitized from Fig.~4
of~\fref{f-g23}; \src{glasser2023_inputs.py}. Section~\ref{sec-age}.
\item\label{f-vd} \src{cdc/variant_days.csv}: the days of the variants;
\src{cdc_variants.py} from~\fref{f-var}. Section~\ref{sec-var}.
\item\label{f-x0} \src{cdc/state_t251_mock.csv}: the initial state (mock);
\src{sero_initial_51x16.py} from~\fref{f-sero}, \fref{f-N}, \fref{f-t2}. CSV without a
header, $13m=663$ rows and $n=16$ columns: 13 blocks of $m$ rows, the compartments in the
order $S$, $P$, $E_s$, $E_p$, $V$, $I_a$, $I_{ps}$, $I_{pm}$, $I_s$, $I_m$, $I_h$, $R$, $M$;
within a block the locations in FIPS order; the columns are the age groups.
Section~\ref{sec-mock}.
\item\label{f-doses} \src{mock/doses_51x16x16_mock.xlsx}: the weekly doses (mock),
sheet \src{by_state} by location and sheet \src{by_state_age} by location and age group;
\src{mock_doses_51x16x16.py} from~\fref{f-N}. Section~\ref{sec-mock}.
\end{enumerate}

\subsection{Time}\label{sec-time}

The time $t$ is the number of days from 31 December 2019, and the model is
started at $t=251$, 7 September 2020~\cite[Section~3]{2023Glasser_Feng}.
Vaccination began on 14 December 2020~\cite[Section~2.2]{2026Feng_Day_Glasser},
$t=349$.

\subsection{Population}\label{sec-pop}

$N_{ij}$ is the resident population of location $i$ and age group $j$ on July 1,
2020, of the Vintage 2024 estimates~\cite{2024Census_scest2024}
(\fref{f-sc}): column \src{POPESTIMATE2020}, both sexes, all races and origins,
summed over single years of age (85+ into 75+). Total: $331\,577\,720$.
(\fref{f-N})

\subsection{Contacts}\label{sec-contacts}

$U$ is the projected contact matrix of~\cite{2017Prem_Cook_Jit} for the United
States, all locations (\fref{f-premx}, \src{MUestimates_all_locations_2.xlsx},
sheet ``United States of America''), balanced as
in~\cite[Section~2.4]{2023Glasser_Feng} by $U_{jq}=\sqrt{M_{jq}M_{qj}}$. Row sums:
$2.2$--$22.1$ contacts per day. (\fref{f-U})

\subsection{Locations}\label{sec-loc}

\subsubsection{Distances}\label{sec-dist}
$d_{ip}$ is the distance in miles between the 2010
centers of population~\cite{2010Census_cenpop} (\fref{f-cenpop}) of locations $i$
and $p$, as in~\cite[Fig.~6]{2023Glasser_Feng}: the geodesic distance on the
WGS84 ellipsoid~\cite{1975Vincenty}. (\fref{f-d})

\subsubsection{Activity factors}\label{sec-L}
$L_i=\ln\rho_i/\ln\rho$~\cite[Section~2.4]{2023Glasser_Feng},
densities in persons per square mile of land. As
in~\cite[Supplement, Section~2a]{2026Feng_Day_Glasser}, $\rho_i$ is the
population-weighted average of the densities of the counties of location $i$,
and $\rho$ the population-weighted average of the $\rho_i$; county populations
on July 1, 2020 (\fref{f-co}), land areas (\fref{f-gaz}). Then $\rho=2253$ and
$L_i\in[0.33,1.27]$. The supplement can also be read with $\rho$ the
population-weighted average of the state densities population$/$area, which
gives larger $L_i$. (\fref{f-L})

\subsubsection{Decay with distance}\label{sec-b}
$b_q$ solves
$\sum_lU_{ql}=\frac1m\sum_{r,k}e^{-b_qd_{rk}}$~\cite[Section~2.2.2.2]{2017Feng_Hill_Curns_Glasser};
computed from $U$ and $d$, $b_q\in[0.0009,0.0087]$ per mile.

\subsubsection{Contact matrix}\label{sec-C}
The mixing matrix $C=(c_{ij,pq})$, of order $mn=816$: $c_{ij,pq}$ is the fraction of
the contacts of a person of location $i$ and age group $j$ that are with persons of
location $p$ and age group $q$, so that every row of $C$ adds up to $1$. How many
contacts a person has is given separately, by the contact rates
$a_{ij}=L_i\sum_lU_{lj}$.

\emph{Source.} The mixing matrix is defined in~\cite[Section~2.4, pp.~5--6]{2023Glasser_Feng}.
For age groups, $c_{jq}=U_{jq}/a_j$ with $a_j=\sum_kU_{jk}$, so that every row adds up to
$1$, the condition $\sum_qc_{jq}=1$ on mixing matrices stated
in~\cite[Section~2.1]{2017Feng_Hill_Curns_Glasser}. For locations alone, the contacts are
$C_{ip}=a_i\,e^{-bd_{ip}}/\sum_ke^{-bd_{ik}}$ and the mixing is
$c_{ip}=C_{ip}/\sum_kC_{ik}$. For locations and age groups,
\cite[Section~2.4, p.~6]{2023Glasser_Feng} gives
\begin{equation}\label{contacts-glasser}
c_{l_ia_j,\,l_pa_q}=\frac{c^{(p)}_{a_ja_q}\,e^{-b_{a_q}d_{l_il_p}}}
{\sum_{r=1}^{n}\sum_{s=1}^{m}c^{(s)}_{a_ja_r}\,e^{-b_{a_r}d_{l_il_s}}},
\qquad
c^{(p)}_{a_ja_q}=c_{a_ja_q}\,\frac{\ln\left(\rho^{(p)}\right)}{\ln\left(\rho\right)}:
\end{equation}
the age-specific mixing within location $p$ is multiplied by the ratio of the natural
logarithms of the location and national population densities, our $L_p$
(Section~\ref{sec-L}), and between locations the contacts decline exponentially with
distance at the rate of the age group met, computed as
in~\cite{2017Feng_Hill_Curns_Glasser}, our $b_q$ (Section~\ref{sec-b}).
By~\cite[Supplement, Section~2a]{2026Feng_Day_Glasser}, the mixing of the model
of~\cite{2026Feng_Day_Glasser} is the mixing used in~\cite{2023Glasser_Feng}, except that
the densities $\rho_i$ are computed from the county densities (Section~\ref{sec-L}).

\emph{The formula used.} In our notation, with $l_i\to i$, $a_j\to j$,
\eqref{contacts-glasser} is
\begin{equation}\label{contacts}
c_{ij,pq}=\frac{L_pU_{jq}e^{-b_qd_{ip}}}{\sum_{p',q'}L_{p'}U_{jq'}e^{-b_{q'}d_{ip'}}}:
\end{equation}
we use the contacts $U_{jq}$ in place of $c_{jq}=U_{jq}/a_j$; the factor $1/a_j$ depends
only on the row $(i,j)$ and cancels between the numerator and the denominator, so
\eqref{contacts} and~\eqref{contacts-glasser} give the same matrix. A person of $(i,j)$
meets age group $q$ in proportion to $U_{jq}$, location $p$ in proportion to its activity
$L_p$, and the contacts decline with the distance $d_{ip}$ at the rate $b_q$ of the age
group met. For fixed $(i,j)$ the denominator is the sum of the numerator over all $(p,q)$,
so every row adds up to $1$. For $n=1$ and $L_p=1$, \eqref{contacts} is the $c_{ip}$
given for locations alone.

\emph{The file.} $C$ is formed from $U$, $L$, $d$ and $b$, and from nothing else, by
\src{build_C_51x16.py} (\fref{f-C}), and \src{seir7_ode.py} loads it from this file.
Its entries are in $[1.5\cdot10^{-23},\,0.34]$, every row adds up to $1$ (to
$3\cdot10^{-15}$), and $C\bm x$ agrees with its computation without forming $C$
(appendix of \src{stratified_seir_ngm.tex}) to $2\cdot10^{-15}$.
The rank and the nonzero entries of $C$, and of $C$ after dropping its entries below a
threshold ($\epsilon=2^{-52}$; the rank is the number of singular values, computed by the
SVD, larger than $816\,\epsilon\,\sigma_{\max}$):
\begin{center}
\small
\begin{tabular}{lcc}
 & rank & nonzeros\\\hline
$C$ & 816 & full\\
drop $|c_{ij,pq}|<16\epsilon$ & 816 & $99.62\%$\\
drop $|c_{ij,pq}|<\sqrt{\epsilon}$ & 816 & $95.86\%$
\end{tabular}
\end{center}

\subsubsection{Stringency}\label{sec-s}
$s_i(t)=1-\mathrm{SI}_i(t)/100$, the complement of the
Oxford stringency index of location $i$~\cite[Section~4.1]{2023Glasser_Feng}:
the daily \src{StringencyIndex_Average} of the US states (\src{STATE_TOTAL})
in~\fref{f-ox}, 1 January 2020 -- 31 December 2022, linear between days and
constant after the last. At $t=251$, $s_i\in[0.20,0.62]$. (\fref{f-s})

\subsection{Age-specific values}\label{sec-age}

From Table~S2 of~\cite{2026Feng_Day_Glasser}, which is Table~2
of~\cite{2023Glasser_Feng}: the proportion with mild if any symptoms $p_j$ and
the proportion hospitalized if infected $d_j$~\cite{2020Clark_Jit}, the infection
fatality ratio $\mathrm{IFR}_j$~\cite{2020Levin_Hanage}, and the compliance $c_j^0$,
estimated in~\cite[Section~4.1]{2023Glasser_Feng} from~\cite{2020Jones_YouGov}
(\fref{f-t2}). From these, $q_j=1-(1-p_j)/2$ \cite[Table~2]{2026Feng_Day_Glasser};
$\delta_{s,j}=\gamma_{s\to h}d_j$ and $\delta_{m,j}=\delta_{s,j}/2$ (Table~S1);
$\mathrm{SFR}_j=\mathrm{IFR}_j(1-p_j)$ (Supplement, Section~2); and
$\theta_{h,j}=1-0.7\,\mathrm{SFR}_j$, $\theta_{s,j}=1-0.3\,\mathrm{SFR}_j$, from the
$70\%$ of deaths among inpatients in~\cite[Section~3]{2023Glasser_Feng}.

The probabilities of infection on contact $\beta_j^0$ are the values for 26 November
2020 shown in~\cite[Fig.~4]{2023Glasser_Feng}, used there for transmission
modeling~\cite[Section~2.3]{2023Glasser_Feng}. The figure gives no table; we
digitized the curve at the 16 age groups (axis calibrated on the tick marks,
values read at the centers of the markers; resolution $5.7\cdot10^{-4}$, accuracy
about $\pm0.003$) (\fref{f-beta}):
\begin{center}
\small
\begin{tabular}{l*{8}{c}}
age & 0--4 & 5--9 & 10--14 & 15--19 & 20--24 & 25--29 & 30--34 & 35--39\\
$\beta_j^0$ & 0.121 & 0.078 & 0.064 & 0.054 & 0.080 & 0.071 & 0.073 & 0.067\\\hline
age & 40--44 & 45--49 & 50--54 & 55--59 & 60--64 & 65--69 & 70--74 & 75+\\
$\beta_j^0$ & 0.067 & 0.078 & 0.078 & 0.095 & 0.142 & 0.170 & 0.187 & 0.333
\end{tabular}
\end{center}

\subsection{Constants}\label{sec-const}

From the supplement of~\cite{2026Feng_Day_Glasser}, per day:
$\xi=\gamma_{s\to h}=1/2$, $\gamma_a=1/7$, $\gamma_s=\gamma_m=1/5$, $\gamma_h=1/10$,
$\omega=\omega_v=1/365$, $b_S=b_I=0.31$ (Table~S1);
$\eta_a,\eta_{ps},\eta_{pm},\eta_s,\eta_m,\eta_h=0.5,1.25,1,1,0.75,0.1$ (Section~2).

\subsection{Variants}\label{sec-var}

The values that change with the dominant variant are set as in Table~S4
of~\cite{2026Feng_Day_Glasser} (a blank there keeps the previous value; $k$ for
Wuhan from Table~S4, where Table~S1 gives $1/4$):
\begin{center}
\small
\begin{tabular}{lccccc}
 & Wuhan & Alpha & Delta & BS1 & BA5\\\hline
factor of $\beta_j^0$ & 1 & 1.25 & 1.5 & 2 & 1\\
$\varepsilon_1=\varepsilon_2=\varepsilon_3$ & 0.9 & 0.8 & 0.7 & 0.6 & 0.5\\
factor of $c_j^0$ & 1 & 0.8 & 0.6 & 0.4 & 0.2\\
$1/k$ (days) & 5 & 5 & 4 & 4 & 3\\
factor of $d_j$ and $1-\theta_{w,j}$ & 1 & 1 & 1 & 0.67 & 0.67\\\hline
from (date) & & 3 Apr 2021 & 3 Jul 2021 & 25 Dec 2021 & 2 Jul 2022\\
from ($t$) & & 459 & 550 & 725 & 914
\end{tabular}
\end{center}
A variant is taken to apply from the end of the first week in which its share of
the US sequences exceeds $1/2$, in the weighted estimates of~\fref{f-var} (for
each week the estimates of the latest creation date; lineages, named as in the CDC data,
grouped as Alpha: B.1.1.7, Q; Delta: B.1.617.2, AY; BS1: B.1.1.529, BA.1; BA5: BA.5, BQ,
BF, BE). The names Wuhan, Alpha, Delta, BS1, BA5 are the column names of Table~S4; the
table gives no lineages, and the lineages assigned to each name and the threshold $1/2$ are
our choice. (\fref{f-vd})

\subsection{Mock data}\label{sec-mock}

\subsubsection{Initial state (mock)}\label{sec-x0}
The state at $t=251$. The data: the seroprevalence of the
CDC Nationwide Commercial Laboratory Seroprevalence
Survey~\cite{2021Bajema_Wiegand,2023CDC_sero} (\fref{f-sero}), rounds 3
(26 August -- 9 September 2020) and 4 (9 -- 18 September 2020), by location and
the 4 age groups 0--17, 18--49, 50--64, 65+. The mock rules, in place of the
procedure of~\cite[Section~3]{2023Glasser_Feng}: (1) from 4 to 16 age groups,
each 5-year group takes the value of the survey group that contains it, and
15--19 takes $3/5$ of 0--17 and $2/5$ of 18--49; a missing or suppressed value is
replaced by the national one of the same round and age group; (2) the
seropositive fraction $r_{ij}$ is that of round 3, and
$\lambda_{ij}=\max\left\{-\ln\left((1-r^{(4)}_{ij})/(1-r^{(3)}_{ij})\right)/11.5,\ 10^{-4}\right\}$
per day; (3) $R=rN$, and $E_s=\lambda S/k$, $I_a=kpE_s/\gamma_a$,
$I_{ps}=k(1-p)E_s/\xi$, $I_s=\xi I_{ps}/(\gamma_s+\delta_s)$ (the quasi-steady
form of the 12-day sums of~\cite[Section~3]{2023Glasser_Feng}), with $S$ such
that the 13 compartments add up to $N_{ij}$, and the other compartments $0$.
This gives $6.6\%$ seropositive. (\fref{f-x0})

\subsubsection{Weekly doses (mock)}\label{sec-doses}
Made up, by location, age group and dose, for
the 126 weeks from $t=349$ through the end of the Public Health Emergency,
11 May 2023~\cite[Section~2.2, p.~4]{2026Feng_Day_Glasser} (\fref{f-doses}): age group
$j$ receives $\mathrm{cov}_jN_{ij}$ first doses evenly over 12 weeks from its
start week ($\mathrm{cov}=0.9$ from week 1 for 65+; $0.8$ from week 9 for 50--64;
$0.7$ from week 15 for 20--49; $0.55$, $0.45$, $0.3$, $0.1$ from weeks 17, 22, 47,
79 for 15--19, 10--14, 5--9, 0--4), and after these 12 weeks $0.002\,N_{ij}$ first
doses per week until the end, so that every location has first doses in every
week; second doses $0.9$ times the first doses four
weeks earlier; third doses half of all second doses of the group, evenly over 20
weeks from week 40 (65+) or 48 (20--64); each location, dose and week multiplied
by $e^{0.1z}$, $z$ standard normal; then each $D_{k,ij}$ rounded to the nearest
integer. To be replaced by the doses reported to the
CDC~\cite[Section~2.2]{2026Feng_Day_Glasser}.

The file~\fref{f-doses} has two sheets, each with two header rows: in the first,
``Week 1'', \dots, ``Week 126'', each over three columns (week $w$ starts on day
$t=349+7(w-1)$); in the second, ``Dose 1'', ``Dose 2'', ``Dose 3+'' under each
week. Sheet 1, \src{by_state}: the column ``State'' and one row for each of the
51 locations, the doses $D_{k,i}=\sum_jD_{k,ij}$ of all ages. Sheet 2,
\src{by_state_age}: the columns ``State'', ``Age'' and one row for each of the
816 groups $(i,j)$, the doses $D_{k,ij}$. All entries are integers, stored as
floating point numbers. Sheet 1 is the sum of sheet 2 over the age groups, exactly
(checked by \src{mock_doses_51x16x16.py} after the file is written).
The optimal strategy uses only sheet 1; sheet 2 is used only by the actual
strategy (Remark~\ref{rem-actual}).

\section{The model and the computation}\label{sec-model-comp}

The program and its files, in \src{MathBio/code/python/}. The program reads the
built files \fref{f-N}--\fref{f-doses} of Section~\ref{sec-built}.
\begin{enumerate}[label=F2-\arabic*, ref=F2-\arabic*, leftmargin=*]
\item\label{f-prog} \src{seir7_ode.py}: the program, run with the arguments
\texttt{simulate optimal} (Section~\ref{sec-comp}).
\item\label{f-alt} \src{alternating.py}: the optimization of each week, called
by~\fref{f-prog}.
\item\label{f-out-o} \src{data/simulate_weekly_51x16_optimal.csv}: the national
totals of the 13 compartments at the start of each week.
\item\label{f-out-rho} \src{data/simulate_weekly_51x16_optimal_rho.csv}: the
spectral radius of each week, at the uniform allocation and at the solution.
\item\label{f-ckpt} \src{intermediate/simulate_optimal_checkpoint.npz}: after every
week, the state, the week and the results so far; with the argument
\texttt{resume} the computation continues from the last completed week.
\end{enumerate}
From \src{MathBio/code/python/}, the computation is run, and after an interruption
resumed, by
\begin{center}
\small
\texttt{python3 -u seir7\_ode.py simulate optimal 2>\&1 | tee intermediate/seir7\_ode\_simulate\_optimal.output}\\
\texttt{python3 -u seir7\_ode.py simulate optimal resume 2>\&1 | tee -a intermediate/seir7\_ode\_simulate\_optimal.output}
\end{center}
The inputs \fref{f-N}--\fref{f-doses} are rebuilt from the downloads by the scripts
named in Section~\ref{sec-built}.

\subsection{The model}\label{sec-model}

The system of~\cite[Section~2.1 and Fig.~1]{2026Feng_Day_Glasser}, with the
location $i$ and the age group $j$ written separately. Each group $(i,j)$ has
13 compartments: susceptible $S$; vaccinated once $V$; partially immune $P$;
exposed $E_s$ (naive) and $E_p$ (partially immune); mild $I_a$; pre-severe and
severe $I_{ps}$, $I_s$; pre-moderate and moderate $I_{pm}$, $I_m$; hospitalized
$I_h$; recovered $R$; dead $M$. The force of infection on group $(i,j)$ is
\begin{equation}\label{lambda}
\lambda_{ij}=\psi_{S,j}\,s_i\,a_{ij}\,\beta_j\sum_{p,q}c_{ij,pq}\,\psi_{I,q}\,
\frac{\sum_x\eta_xI_{x,pq}}{N_{pq}},\qquad x\in\{a,ps,pm,s,m,h\},
\end{equation}
with $a_{ij}$ and $C=(c_{ij,pq})$ of Section~\ref{sec-loc}, $s_i(t)$ of
Section~\ref{sec-loc}, and $\psi_{S,j}=1-c_jb_S$, $\psi_{I,j}=1-c_jb_I$,
$c_j=c_j^0$ and $\beta_j=\beta_j^0$ times the factors of the variant at $t$
(Sections~\ref{sec-age} and~\ref{sec-var}). For every group (indices $i,j$
omitted; $\lambda=\lambda_{ij}(t)$, $\varepsilon=\varepsilon(t)$ of the variant,
$\chi_k=\chi_{k,ij}$ the vaccination rates of dose $k$),
\begin{equation}\label{ode}
\begin{aligned}
S'&=\omega P+\omega_vV-(\lambda+\chi_1)S, &
V'&=\chi_1S-\left((1-\varepsilon)\lambda+\omega_v+\varepsilon\chi_2\right)V,\\
P'&=\omega R-(\omega+\varepsilon\chi_3+\lambda)P, &
E_s'&=\lambda S-kE_s,\\
E_p'&=(1-\varepsilon)\lambda V+\lambda P-kE_p, &
I_a'&=kp_jE_s+kq_jE_p-\gamma_aI_a,\\
I_{ps}'&=k(1-p_j)E_s-\xi I_{ps}, &
I_{pm}'&=k(1-q_j)E_p-\xi I_{pm},\\
I_s'&=\xi I_{ps}-(\gamma_s+\delta_{s,j})I_s, &
I_m'&=\xi I_{pm}-(\gamma_m+\delta_{m,j})I_m,\\
I_h'&=\delta_{s,j}I_s+\delta_{m,j}I_m-\gamma_hI_h, &
M'&=(1-\theta_{s,j})\gamma_sI_s+(1-\theta_{h,j})\gamma_hI_h,
\end{aligned}
\end{equation}
\begin{equation*}
R'=\gamma_aI_a+\theta_{s,j}\gamma_sI_s+\gamma_mI_m+\theta_{h,j}\gamma_hI_h
+\varepsilon\chi_2V+\varepsilon\chi_3P-\omega R,
\end{equation*}
$13mn=10\,608$ unknowns, coupled only through $\lambda$. The 13 right hand sides
of a group add up to $0$, so each group keeps its size $N_{ij}$.

\subsection{The computation}\label{sec-comp}

\begin{enumerate}[leftmargin=*]
\item Solve~\eqref{ode} from $t=251$, with the initial state~\fref{f-x0}, to
$t=349$, the start of vaccination, with $\chi=0$.
\item For each week $[\tau,\tau+7]$, $\tau=349,356,\dots$ (126 weeks), with
$S$, $V$, $P$ the state at $\tau$ and $D_{k,i}$ the mock doses of location $i$ in the
week, sheet 1 of~\fref{f-doses} (Section~\ref{sec-doses}):
\begin{enumerate}[label=(\alph*)]
\item compute the vaccination rates $\chi_{k,ij}$ for the week (the same on every
day of the week; they change only from one week to the next): minimize the
spectral radius of the next generation matrix over $\chi$, subject to
$\sum_j\sigma_{k,ij}\chi_{k,ij}=\frac17D_{k,i}$ for every location $i$ and
dose $k$ (the doses of the location, all ages, at most $0.99\sum_j\sigma_{k,ij}/7$),
with $\sigma_1=S$, $\sigma_2=V$, $\sigma_3=P$, only the groups with
$\sigma_{k,ij}\ge1$ (at least one eligible person) receiving dose $k$, and the bound
$0\le\chi_{k,ij}\le\chi_{\max}=-\frac17\ln(0.01)\approx0.658$ per day (at most $99\%$
of a group in a week) as a constraint of the minimization
(\src{stratified_seir_ngm.tex}, problem (1.1), Section~9 and its Section~9.1);
\item solve~\eqref{ode} on $[\tau,\tau+7]$ with these rates, which gives the state
at $\tau+7$;
\item advance to the next week, $\tau\leftarrow\tau+7$.
\end{enumerate}
\end{enumerate}
The system is solved by the Runge--Kutta method of order 8 of Dormand and
Prince~\cite{1980Dormand_Prince} (\texttt{DOP853} of
\texttt{scipy.integrate.solve\_ivp}~\cite{2020Virtanen_scipy}), restarted at the
days at which the variant changes.

\begin{remark}[The actual strategy]\label{rem-actual}
In~\cite{2026Feng_Day_Glasser} the optimal allocation is compared with the
actual one, the doses reported to the CDC by location, age group and dose. In
the actual strategy no optimization is done: in each week, the fraction
$P_{k,ij}=\min\{D_{k,ij}/X_{k,ij}(\tau),0.99\}$ of the eligible people
$X_1=S$, $X_2=V$, $X_3=P$ of group $(i,j)$ is vaccinated, that is,
$\chi_{k,ij}=-\frac17\ln\left(1-P_{k,ij}\right)$, and the system~\eqref{ode} is
solved with these rates. With the doses by age group
of~\cite[Section~2.2]{2026Feng_Day_Glasser}, it gives the reference against which
the optimal allocation is measured (illnesses, hospitalizations and deaths
averted). It is not computed here; only the optimal strategy is run. We have no
reported doses: computed with the mock doses $D_{k,ij}$ of~\fref{f-doses} (sheet 2), the actual
strategy would give a mock reference only.
\end{remark}

\section[The two parameter sets: Unknown 2020/21 and Feng, Day, Glasser 2026; Glasser et al.\ 2023]{The two parameter sets: Unknown 2020/21 and~\cite{2026Feng_Day_Glasser},~\cite{2023Glasser_Feng}}\label{sec-params}

The model is computed with one of two parameter sets: the values of Unknown 2020/21
(set \src{2021}, the default) and the values used in~\cite{2026Feng_Day_Glasser}
and~\cite{2023Glasser_Feng}, as described in Section~\ref{sec-data} (set \src{feng2026}). Both are
in \src{weekly_simulation/input/parameters.xlsx}, one column per set (sheets \src{scalars},
\src{by_age}, \src{over_time}); the set is chosen by \src{params=2021} or
\src{params=feng2026} in \src{seir7_ode.py}, \src{ngm_seir7.py} and
\src{plot_weekly_optimal.py}, and its name ends the names of the output files.

\subsection{Constants of the two sets}\label{sec-params-const}

Per day, unless stated otherwise:
\begin{center}
\small
\begin{tabular}{lcc}
 & Unknown 2020/21 & \cite{2026Feng_Day_Glasser}, \cite{2023Glasser_Feng}\\\hline
$\xi$, $\gamma_{s\to h}$ & $1/2$, $1/2$ & $1/2$, $1/2$\\
$\gamma_a$, $\gamma_s$, $\gamma_m$, $\gamma_h$ & $1/7$, $1/5$, $1/5$, $1/10$ & $1/7$, $1/5$, $1/5$, $1/10$\\
$\omega=\omega_v$ & $1/182.5$ & $1/365$\\
$b_S=b_I$ & 0.31 & 0.31\\
$\eta_a,\eta_{ps},\eta_{pm},\eta_s,\eta_m,\eta_h$ & 0.5, 1.25, 1, 1, 0.75, 0.1 & 0.5, 1.25, 1, 1, 0.75, 0.1\\
deaths of $\mathrm{SFR}_j$: outside, in hospital & 0.33, 0.67 & 0.3, 0.7\\
start of vaccination & $t=347$ (12 Dec 2020) & $t=349$ (14 Dec 2020)
\end{tabular}
\end{center}
That is, $\theta_{s,j}=1-0.33\,\mathrm{SFR}_j$ and $\theta_{h,j}=1-0.67\,\mathrm{SFR}_j$ in the
set of Unknown 2020/21. In both sets $q_j=1-(1-p_j)/2$, $\delta_{s,j}=\gamma_{s\to h}d_j$ and
$\delta_{m,j}=\delta_{s,j}/2$. The doses of week $w$ are given from
$t=t_{\mathrm{vacc}}+7(w-1)$, with $t_{\mathrm{vacc}}$ the start of vaccination of the
set.

\subsection{Age-specific values of the two sets}\label{sec-params-age}

The values $p_j$, $d_j$ and $\mathrm{IFR}_j$ are the same in both sets
(Section~\ref{sec-age}); the compliance $c_j^0$ of Unknown 2020/21 is given with 16 digits and
agrees with the compliance described in Section~\ref{sec-age} to the 9 digits given there.
The probabilities of infection on contact $\beta_j^0$ differ: the values of Unknown 2020/21,
and the values from Fig.~4~\cite{2023Glasser_Feng}, as described in Section~\ref{sec-age}:
\begin{center}
\small
\begin{tabular}{l*{8}{c}}
age & 0--4 & 5--9 & 10--14 & 15--19 & 20--24 & 25--29 & 30--34 & 35--39\\\hline
Unknown 2020/21 & 0.1241 & 0.0815 & 0.0645 & 0.0511 & 0.0772 & 0.0762 & 0.0733 & 0.0683\\
from Fig.~4~\cite{2023Glasser_Feng} & 0.1211 & 0.0779 & 0.0637 & 0.0535 & 0.0796 & 0.0711 & 0.0728 & 0.0671\\\hline
age & 40--44 & 45--49 & 50--54 & 55--59 & 60--64 & 65--69 & 70--74 & 75+\\\hline
Unknown 2020/21 & 0.0689 & 0.0794 & 0.0793 & 0.0951 & 0.1426 & 0.1708 & 0.1881 & 0.3305\\
from Fig.~4~\cite{2023Glasser_Feng} & 0.0671 & 0.0779 & 0.0779 & 0.0950 & 0.1422 & 0.1701 & 0.1871 & 0.3327
\end{tabular}
\end{center}
The values from Fig.~4~\cite{2023Glasser_Feng} differ from the values of Unknown 2020/21 by $-6.7\%$ (25--29) to $+4.7\%$
(15--19), and by less than $1\%$ from age 55.

\subsection{Values that change over time in the two sets}\label{sec-params-var}

In the set of Unknown 2020/21 each value has its own days of change, and the values after
1 July 2021 are kept to the end (there are no values for BS1 and BA5); a value given
from day $t$ holds for the days after $t$:
\begin{center}
\small
\begin{tabular}{lcccc}
 & start & from $t=398$ & from $t=460$ & from $t=548$\\
 & & (1 Feb 2021) & (4 Apr 2021) & (1 Jul 2021)\\\hline
factor of $\beta_j^0$ & 1 & 1.25 & & 1.5\\
$\varepsilon_1=\varepsilon_2=\varepsilon_3$ & 0.9 & 0.8 & & 0.7\\
factor of $c_j^0$ & 1 & & 0.8 & \\
$1/k$ (days) & 5 & & & \\
factor of $d_j$ and $1-\theta_{w,j}$ & 1 & & &
\end{tabular}
\end{center}
A blank keeps the previous value. In the set used in~\cite{2026Feng_Day_Glasser}, as
described in Section~\ref{sec-var}, all values change together on the days of the variants,
$t=459$, $550$, $725$, $914$. The two sets differ in the day of the first change ($t=398$
for $\beta$ and $\varepsilon$, $t=460$ for the compliance, against $t=459$ for all), in
the factor of $c_j^0$ under Delta (0.8 against 0.6), in $k$ under Delta ($1/5$ against
$1/4$), and after Delta.

\bibliographystyle{plainurl}
\bibliography{bib_mb}
\end{document}
